\section{Where the Adapter Switches}\label{sec:switch}

\subsection{Student Outputs}\label{sec:outputs}

We first examine what changes in the student's outputs when the teacher stops thinking.
Each pair's teacher drops \texttt{</think>} at a different $\lambda$, and the reflective markers in its solutions fall from 13--16 to 0.1--3.7 per 1,000 words.
The reflective markers of the students trained on that teacher fall from 14--44 to 0.07--5.1 per 1,000 words (Table~\ref{tab:1}).
The students' reflective markers closely track the amount present in the teacher's solutions: their changes over $\lambda$ agree with the teacher's at $r = 0.82$--$1.00$ in each of the nine cells, and at $r = 0.87$ with all nine cells pooled (Table~\ref{tab:1}).
Over the same range, responses become 2--7 times shorter and truncation falls from 51--87\% to 4--22\%.
Accuracy rises by 1.7--19.2 percentage points (pp), but it does not rise when only completed responses are scored (Table~\ref{tab:2}).

The student's outputs thus change with the teacher's solutions: reflective markers and length decrease together with $\lambda$, and most of the accuracy gain is the share recovered from truncation.
Because the completed responses at $\lambda=0$ are concentrated on easy problems, we do not interpret even the sign of the completed-only values.
The teacher switch is observed in all three pairs, consistent with studies that report this switch \citep{wu2026interpolation,lan2025spectrum}.
In short, the student's reflective markers and length decrease together, following the teacher's solutions, so the outputs alone cannot tell whether the way of thinking changed or only the length did.

\begin{table*}[t]
\centering
\small
% Table 1 (wide: use table*), generated from overleaf/04_where_the_adapter_switches.md
\resizebox{\linewidth}{!}{%
\begin{tabular}{llrrrrrr}
\toprule
Pair & Model & $\lambda$ = 0 & $\lambda$ = 0.2 & $\lambda$ = 0.4 & $\lambda$ = 0.6 & $\lambda$ = 0.8 & $\lambda$ = 1 \\
\midrule
A & Teacher think rate & 0.83 & 0.91 & 0.00 & 0.00 & 0.00 & 0.00 \\
 & Teacher markers & 15.5 & 5.0 & 1.0 & 0.5 & 0.4 & 0.2 \\
 & Qwen2.5-0.5B & 21.8 $\pm$ 2.5 & 5.4 $\pm$ 0.6 & 0.91 $\pm$ 0.66 & 0.45 $\pm$ 0.23 & 0.41 $\pm$ 0.36 & 0.48 $\pm$ 0.61 \\
 & Llama-3.2-1B & 30.7 $\pm$ 1.2 & 6.6 $\pm$ 0.9 & 0.56 $\pm$ 0.08 & 0.62 $\pm$ 0.53 & 0.20 $\pm$ 0.06 & 0.12 $\pm$ 0.03 \\
 & OLMo-2-1B & 15.0 $\pm$ 1.5 & 6.1 $\pm$ 1.1 & 0.74 $\pm$ 0.08 & 0.60 $\pm$ 0.35 & 0.24 $\pm$ 0.12 & 0.41 $\pm$ 0.13 \\
\midrule
B & Teacher think rate & 0.73 & 0.83 & 0.91 & 0.01 & 0.00 & 0.00 \\
 & Teacher markers & 13.3 & 12.3 & 10.9 & 5.0 & 4.2 & 3.7 \\
 & Qwen2.5-0.5B & 19.2 $\pm$ 2.2 & 12.8 $\pm$ 2.4 & 12.1 $\pm$ 2.7 & 6.8 $\pm$ 4.5 & 6.3 $\pm$ 2.9 & 4.7 $\pm$ 2.1 \\
 & Llama-3.2-1B & 44.4 $\pm$ 26.0 & 15.1 $\pm$ 3.5 & 17.1 $\pm$ 1.9 & 4.0 $\pm$ 0.5 & 5.9 $\pm$ 0.7 & 2.8 $\pm$ 0.7 \\
 & OLMo-2-1B & 16.8 $\pm$ 2.4 & 9.9 $\pm$ 0.9 & 10.4 $\pm$ 0.7 & 4.9 $\pm$ 0.4 & 6.1 $\pm$ 0.7 & 5.1 $\pm$ 2.0 \\
\midrule
C & Teacher think rate & 0.90 & 0.90 & 0.89 & 0.27 & 0.00 & 0.00 \\
 & Teacher markers & 14.2 & 12.6 & 11.2 & 5.0 & 0.6 & 0.1 \\
 & Qwen2.5-0.5B & 16.6 $\pm$ 1.7 & 15.5 $\pm$ 2.1 & 16.0 $\pm$ 0.8 & 9.7 $\pm$ 1.3 & 0.17 $\pm$ 0.20 & 0.07 $\pm$ 0.06 \\
 & Llama-3.2-1B & 19.1 $\pm$ 1.1 & 17.8 $\pm$ 1.7 & 20.2 $\pm$ 2.4 & 13.6 $\pm$ 0.4 & 0.58 $\pm$ 0.40 & 0.14 $\pm$ 0.14 \\
 & OLMo-2-1B & 14.4 $\pm$ 1.5 & 13.8 $\pm$ 1.0 & 12.9 $\pm$ 0.7 & 7.4 $\pm$ 1.8 & 0.39 $\pm$ 0.16 & 0.33 $\pm$ 0.22 \\
\bottomrule
\end{tabular}}

\caption{Teacher think rate (share of responses containing \texttt{</think>}) and frequency of reflective markers in the teacher's solutions for each pair, shown alongside the reflective-marker frequency (per 1,000 words) of the students trained on that teacher, across $\lambda$. The students' reflective markers decrease following the frequency in the teacher's solutions. In pair A, the teacher's reflective markers decrease one interval earlier than the tag. Values after $\pm$ are standard deviations across three seeds.}
\label{tab:1}
\end{table*}

\begin{table*}[t]
\centering
\small
\resizebox{\linewidth}{!}{% Table 2 (fits one column or use table*), generated from overleaf/04_where_the_adapter_switches.md
\begin{tabular}{llrrrr}
\toprule
Pair & Student & Tokens ($\lambda$ = 0 $\rightarrow$ $\lambda$ = 1) & Truncation rate ($\lambda$ = 0 $\rightarrow$ $\lambda$ = 1) & $\Delta$ Acc. (pp) & $\Delta$ Acc., completed only (pp) \\
\midrule
A & Qwen2.5-0.5B & 5,916 $\pm$ 214 $\rightarrow$ 811 $\pm$ 64 & 0.70 $\pm$ 0.04 $\rightarrow$ 0.04 $\pm$ 0.01 & +19.2 $\pm$ 0.9 & $-$4.4 $\pm$ 1.0 \\
 & Llama-3.2-1B & 6,969 $\pm$ 77 $\rightarrow$ 944 $\pm$ 72 & 0.87 $\pm$ 0.02 $\rightarrow$ 0.05 $\pm$ 0.01 & +18.1 $\pm$ 3.2 & $-$14.3 $\pm$ 2.8 \\
 & OLMo-2-1B & 3,015 $\pm$ 86 $\rightarrow$ 683 $\pm$ 23 & 0.61 $\pm$ 0.04 $\rightarrow$ 0.05 $\pm$ 0.01 & +9.6 $\pm$ 1.5 & +1.4 $\pm$ 3.5 \\
\midrule
B & Qwen2.5-0.5B & 5,226 $\pm$ 82 $\rightarrow$ 2,302 $\pm$ 94 & 0.55 $\pm$ 0.01 $\rightarrow$ 0.22 $\pm$ 0.01 & +8.8 $\pm$ 1.1 & $-$2.8 $\pm$ 1.6 \\
 & Llama-3.2-1B & 5,095 $\pm$ 164 $\rightarrow$ 1,765 $\pm$ 78 & 0.52 $\pm$ 0.04 $\rightarrow$ 0.14 $\pm$ 0.01 & +6.5 $\pm$ 1.7 & $-$6.2 $\pm$ 2.5 \\
 & OLMo-2-1B & 2,827 $\pm$ 25 $\rightarrow$ 1,207 $\pm$ 40 & 0.51 $\pm$ 0.03 $\rightarrow$ 0.16 $\pm$ 0.01 & +5.7 $\pm$ 2.1 & $-$0.2 $\pm$ 2.4 \\
\midrule
C & Qwen2.5-0.5B & 5,231 $\pm$ 141 $\rightarrow$ 1,399 $\pm$ 156 & 0.59 $\pm$ 0.02 $\rightarrow$ 0.13 $\pm$ 0.02 & +8.7 $\pm$ 1.4 & $-$5.8 $\pm$ 1.9 \\
 & Llama-3.2-1B & 6,149 $\pm$ 103 $\rightarrow$ 939 $\pm$ 96 & 0.73 $\pm$ 0.02 $\rightarrow$ 0.06 $\pm$ 0.01 & +15.3 $\pm$ 1.6 & $-$10.2 $\pm$ 3.2 \\
 & OLMo-2-1B & 2,949 $\pm$ 50 $\rightarrow$ 910 $\pm$ 37 & 0.58 $\pm$ 0.02 $\rightarrow$ 0.14 $\pm$ 0.01 & +1.7 $\pm$ 3.0 & $-$4.6 $\pm$ 4.9 \\
\bottomrule
\end{tabular}
}
\caption{MATH-500 response length, truncation rate, and accuracy change (pp) of the $\lambda=0$ and $\lambda=1$ students. The accuracy change is reported both over all responses and with only completed responses scored. The release from truncation as responses become shorter alone explains most of the accuracy gain. Values after $\pm$ are standard deviations across three seeds.}
\label{tab:2}
\end{table*}

\subsection{Adapter Directions}\label{sec:directions}

We therefore look at the adapter left by training instead of the outputs, and ask whether the adapter sharply changes the directions along which it updates the weights at the $\lambda$ where the teacher stops thinking.
$R$ drops most at $\lambda = 0.2 \to 0.4$ in pair A and at $0.4 \to 0.6$ in pair B, while in pair C it drops over two steps, $0.4 \to 0.6$ and $0.6 \to 0.8$, as the teacher does (Table~\ref{tab:3}).
In the eight cells other than A/Qwen, the per-interval drops agree with the drops in the teacher's think rate at $r = 0.80$--$0.95$, and the exact test pooling the three pairs gives $p = 0.0035$; with only five $\lambda$ intervals, a single pair reaches only $p = 0.05$--$0.20$ (Table~\ref{tab:4}).
Measuring the teacher by the drop in reflective markers in its solutions instead of the think rate also gives agreement, with $p = 0.0032$ pooled over the three pairs.

Since the $\lambda$ at which $R$ drops sharply differs across pairs and each pair follows its own teacher, the change in the adapter follows the teacher stopping thinking rather than the value of $\lambda$ itself.
In pair C, both $R$ and the students' reflective markers follow the teacher's two-step collapse.
In pair A, $R$ already drops to 0.75--0.81 at $\lambda=0.2$, where the teacher still uses \texttt{</think>} in 91\% of responses; at this point the reflective markers in the teacher's solutions have also already fallen from 15.5 to 5.0 per 1,000 words (Table~\ref{tab:1}).
Measuring the teacher by reflective markers gives a similar correlation for pair A (mean $r$ $0.72 \to 0.66$), so we cannot tell whether $R$ follows the tag or the reflective markers.
For A/Qwen, $R$ falls only to 0.62 and then rises again.
The numerator (overlap between adapters trained on different teachers) falls from 0.31 to 0.045, as much as in the other cells, but the denominator (overlap between adapters differing only in data) also collapses from 0.39 to 0.07, which masks the ratio.
In short, in all three pairs the adapter sharply changes the directions along which it updates the weights at the $\lambda$ where the teacher stops thinking.
However, because $R$ measures the overall overlap as a single value, it cannot tell whether all 16 directions changed wholesale or only some of them.

\begin{table*}[t]
\centering
\small
% Table 3 (wide: use table*), generated from overleaf/04_where_the_adapter_switches.md
\resizebox{\linewidth}{!}{%
\begin{tabular}{llrrrrrr}
\toprule
Pair & Model & $\lambda$ = 0 & $\lambda$ = 0.2 & $\lambda$ = 0.4 & $\lambda$ = 0.6 & $\lambda$ = 0.8 & $\lambda$ = 1 \\
\midrule
A & Teacher think rate & 0.83 & 0.91 & 0.00 & 0.00 & 0.00 & 0.00 \\
 & Qwen2.5-0.5B & 1.00 & 0.81 $\pm$ 0.002 & 0.62 $\pm$ 0.01 & 0.71 $\pm$ 0.007 & 0.74 $\pm$ 0.005 & 0.74 $\pm$ 0.007 \\
 & Llama-3.2-1B & 1.00 & 0.75 $\pm$ 0.004 & 0.26 $\pm$ 0.004 & 0.20 $\pm$ 0.004 & 0.18 $\pm$ 0.001 & 0.17 $\pm$ 0.002 \\
 & OLMo-2-1B & 1.00 & 0.77 $\pm$ 0.008 & 0.38 $\pm$ 0.005 & 0.31 $\pm$ 0.003 & 0.28 $\pm$ 0.004 & 0.27 $\pm$ 0.007 \\
\midrule
B & Teacher think rate & 0.73 & 0.83 & 0.91 & 0.01 & 0.00 & 0.00 \\
 & Qwen2.5-0.5B & 1.00 & 0.97 $\pm$ 0.005 & 0.85 $\pm$ 0.007 & 0.42 $\pm$ 0.008 & 0.25 $\pm$ 0.004 & 0.23 $\pm$ 0.003 \\
 & Llama-3.2-1B & 1.00 & 0.97 $\pm$ 0.003 & 0.84 $\pm$ 0.004 & 0.35 $\pm$ 0.003 & 0.19 $\pm$ 0.001 & 0.18 $\pm$ 0.001 \\
 & OLMo-2-1B & 1.00 & 0.97 $\pm$ 0.004 & 0.85 $\pm$ 0.005 & 0.38 $\pm$ 0.008 & 0.21 $\pm$ 0.002 & 0.19 $\pm$ 0.004 \\
\midrule
C & Teacher think rate & 0.90 & 0.90 & 0.89 & 0.27 & 0.00 & 0.00 \\
 & Qwen2.5-0.5B & 1.00 & 0.98 $\pm$ 0.002 & 0.89 $\pm$ 0.008 & 0.55 $\pm$ 0.009 & 0.26 $\pm$ 0.01 & 0.18 $\pm$ 0.005 \\
 & Llama-3.2-1B & 1.00 & 0.96 $\pm$ 0.004 & 0.84 $\pm$ 0.009 & 0.55 $\pm$ 0.002 & 0.24 $\pm$ 0.01 & 0.16 $\pm$ 0.008 \\
 & OLMo-2-1B & 1.00 & 0.97 $\pm$ 0.003 & 0.86 $\pm$ 0.02 & 0.58 $\pm$ 0.003 & 0.26 $\pm$ 0.004 & 0.18 $\pm$ 0.001 \\
\bottomrule
\end{tabular}}

\caption{Teacher think rate and $R(\lambda)$ of the student adapters for each pair, side by side. $R$ measures how much an adapter uses the same directions as the $\lambda=0$ adapter, relative to a baseline that changes only the data; 1 means the directions are unchanged, and values near 0 mean the directions changed sharply. Each pair's $R$ drops most in the interval where its own teacher drops the tag. Values after $\pm$ are standard deviations across three seeds.}
\label{tab:3}
\end{table*}

\begin{table}[t]
\centering
\small
\resizebox{\columnwidth}{!}{% Table 4 (fits one column or use table*), generated from overleaf/04_where_the_adapter_switches.md
\begin{tabular}{lrrrrr}
\toprule
Pair & Qwen2.5-0.5B & Llama-3.2-1B & OLMo-2-1B & Mean & p (permutation) \\
\midrule
A & 0.53 & 0.84 & 0.80 & 0.72 & 0.20 \\
B & 0.94 & 0.95 & 0.94 & 0.94 & 0.05 \\
C & 0.91 & 0.83 & 0.81 & 0.85 & 0.07 \\
\midrule
Pooled &  &  &  & 0.84 & 0.0035 \\
\bottomrule
\end{tabular}
}
\caption{Correlation $r$ between the drop in teacher think rate and the drop in $R$ in each $\lambda$ interval, and the $p$ of the permutation test. With only five intervals, no single pair is significant, but pooling the three pairs gives $p = 0.0035$, showing that the adapter switch occurs in the same $\lambda$ interval as the teacher switch.}
\label{tab:4}
\end{table}

\subsection{Directions Shared Before the Switch}\label{sec:shared}

We therefore examine which of the 16 directions the pre-switch adapters were using in common.
The pre-switch adapters always share 5.4--7.7 of the 16 directions.
Pre-switch adapters use 41--57\% of this shared region, whereas post-switch adapters use only 5--19\%, against a chance level of 0.5--1.8\% (Table~\ref{tab:5}).
Conversely, when the region is built from post-switch adapters, it spans only 1.3--3.3 dimensions in pairs A and C, and only pair B forms 4.2--5.5 dimensions.

Two adapters with different seeds overlap in only 3--4 directions, yet pooling many adapters yields 5--8 dimensions; thus, although individual directions fluctuate, the region the adapters use in common is fixed \citep{gerasimov2026unstable}.
This agrees with the report that long solutions train the student along more concentrated directions \citep{lan2026midistillation}; here we further show that these directions form the same region across many adapters and that this region empties at the $\lambda$ where the teacher stops thinking.
After the switch, no common region of comparable size newly forms.
The 4--5 dimensions in pair B may result from the short-answer teacher teaching a consistent solution style, but we have not verified this.
In short, the shared region that the pre-switch adapters used together empties after the switch, but having moved together alone does not tell whether this region carries the way of thinking or the length.

\begin{table*}[t]
\centering
\small
% Table 5 (wide: use table*), generated from overleaf/04_where_the_adapter_switches.md
\resizebox{\linewidth}{!}{%
\begin{tabular}{llrrrrrr}
\toprule
Pair & Student & Transition $\lambda$ & Shared dims (pre) & Coverage: pre (LOO) & Coverage: post & Coverage: chance & Shared dims (post) \\
\midrule
A & Qwen2.5-0.5B & 0.4 & 5.4 & 0.41 $\pm$ 0.003 & 0.05 $\pm$ 0.0002 & 0.018 & 1.3 \\
 & Llama-3.2-1B & 0.4 & 6.2 & 0.43 $\pm$ 0.004 & 0.08 $\pm$ 0.001 & 0.007 & 3.3 \\
 & OLMo-2-1B & 0.4 & 6.3 & 0.46 $\pm$ 0.004 & 0.09 $\pm$ 0.0006 & 0.005 & 2.4 \\
\midrule
B & Qwen2.5-0.5B & 0.6 & 6.9 & 0.52 $\pm$ 0.002 & 0.16 $\pm$ 0.0007 & 0.016 & 4.2 \\
 & Llama-3.2-1B & 0.6 & 7.2 & 0.53 $\pm$ 0.001 & 0.16 $\pm$ 0.0004 & 0.007 & 5.0 \\
 & OLMo-2-1B & 0.6 & 7.7 & 0.55 $\pm$ 0.003 & 0.18 $\pm$ 0.0006 & 0.005 & 5.5 \\
\midrule
C & Qwen2.5-0.5B & 0.6 & 6.7 & 0.53 $\pm$ 0.003 & 0.14 $\pm$ 0.002 & 0.018 & 1.6 \\
 & Llama-3.2-1B & 0.6 & 7.3 & 0.54 $\pm$ 0.002 & 0.14 $\pm$ 0.002 & 0.007 & 1.3 \\
 & OLMo-2-1B & 0.6 & 7.5 & 0.57 $\pm$ 0.003 & 0.19 $\pm$ 0.002 & 0.005 & 3.0 \\
\bottomrule
\end{tabular}}

\caption{Directions shared by the pre-switch adapters, and the extent to which each adapter uses that region. Transition $\lambda$ is the first $\lambda$ after the teacher's think rate first drops by more than half; Shared dims (pre) is the number of shared directions out of 16; Coverage is the fraction of this region covered by pre-switch, post-switch, and random adapters; and Shared dims (post) is, conversely, the size of the region built from post-switch adapters. The coverage of a pre-switch adapter is measured with the region built without that adapter (LOO, leave-one-out). Pre-switch adapters use nearly half of the region, which falls to one fifth or less after the switch, and no common region of comparable size newly forms after the switch. Values after $\pm$ are standard deviations across three seeds.}
\label{tab:5}
\end{table*}
